\section{What Users Experience}\label{sec:qoe}

\subsection{Loss and jitter outweigh delay}
Recent user studies agree on one ordering: loss, jitter and stalls hurt more than tens of milliseconds of
added delay. With Nvidia CloudXR and a first-person shooter (N=30), QoE fell at a round-trip time (RTT) of
75~ms or more, or at loss above 6\%; combined, 25~ms with 4\% loss was already harmful \cite{W4402156718}.
In a social XR application (N=28), RTT up to 77~ms had no significant effect, while loss above 2\% at
higher RTT did, and even jitter of one standard deviation degraded QoE at RTTs beyond 52~ms
\cite{W4415822410}. With a Quest~3 streaming over Air Link (N=16), loss harmed experience more than delay
across almost all measured aspects \cite{W4404916373}. In Rec Room and VRChat, the virtual environment
froze under loss but not under latency \cite{W4404103205}.

The ordering follows from how the headset splits latency. Reprojection serves head rotation locally, so
network delay reaches only \emph{interaction latency}: aiming, grabbing, or seeing a peer react. That
budget depends on timing structure, not on RTT alone. In split-rendered gaming, 90 and 120~fps tolerated
up to 6~ms of network delay before responsiveness suffered, and 60~fps was the better choice when delay
was high \cite{W4401943426}. The budget is therefore \emph{contested}: requirement statements still ask for
sub-20~ms end-to-end latency \cite{W7126117938}, while users of cloud and social XR tolerated RTTs near
75~ms. Cybersickness depends on content more than on the network: loss raised sickness in Beat Saber but lowered presence in Cubism \cite{W4404916373}, and at most 21\% of viewers
reported dizziness for point clouds in AR \cite{W4387831705}.

\subsection{Experience is social}
Multi-user sessions add effects no single-user study can show. In collaborative VR (20 participants in
pairs), users told impaired from unimpaired sessions but not latency from bursts, task time and sickness
scores did not change, and the user's role shaped perception \cite{W4392600587}. Impairments also raised synchronisation traffic four- to fivefold through retransmissions. Avatars
expose a fidelity trade-off. With 40 participants on Quest~Pro headsets, a light cartoon avatar gave the
highest QoE under impairment; a humanoid avatar that replicates facial expressions lost more under delay,
yet seeing a partner's face softened the loss \cite{W4415821547}. Standards lag behind: ITU-T leads in
modelling social aspects of QoE, but 3GPP, ITU-T and ETSI efforts remain fragmented and centred on system
metrics \cite{W4415822422}.

\subsection{Predicting experience}
Three families of models map what can be observed to what users would report.
\emph{Opinion-trained models} fit ratings from a lab study: cloud VR gaming models reached a Pearson
correlation of 0.93 \cite{W4400981668}, and regression models for a VR shooter a mean absolute error of 0.14 \cite{W4405935991}. Fed with real traces, the latter predicted a mean opinion
score (MOS) of only 2.5--3.0 over 4G and 5G against 4.2--4.4 over wired access. \emph{Content-quality
metrics} cover new media: colour PSNR correlated with point-cloud ratings at 0.84 \cite{W4387831705}, but
common metrics missed the effect of viewing behaviour on volumetric video \cite{W4386172458}.
\emph{Network-side estimators} infer QoE from what an operator sees, such as encrypted 360-degree traffic
over HTTPS and QUIC in 4G and 5G drive tests \cite{W4384158049} or QoS counters on an O-RAN testbed
\cite{W4389923869}, but are validated against application metrics, not users.

All three families inherit one study design: \emph{static impairments}. Emulators apply a constant RTT,
independent loss or a fixed jitter distribution to 16--40 participants, and real 4G and 5G measurements
enter only afterwards, as inputs to models fitted under constant conditions
\cite{W4405935991,W4415822410}. That thresholds learned this way hold under bursty, correlated wireless loss is \emph{assumed}.
